\documentclass[12pt,a4paper]{article}

\usepackage{fontawesome5}
\providecommand{\faPencilAlt}{\faIcon{pencil-alt}}

\usepackage{styles/couleurs}
\usepackage{styles/mise_en_page}
\usepackage{styles/boites}
\usepackage{styles/maths}
\usepackage{styles/figures}

\usepackage{enumitem}
\usepackage{array}
\usepackage{tabularx}
\usepackage{amsmath}
\usepackage{amssymb}
\usepackage{graphicx}
\usepackage{tikz}
\usepackage{xparse}

\usetikzlibrary{calc,arrows.meta}

\setlist[itemize]{leftmargin=1.4em,itemsep=0.15em,topsep=0.2em}
\setlist[enumerate]{leftmargin=*,itemsep=0.22em,topsep=0.25em}

\newcommand{\lignepointillee}{%
  \noindent\dotfill\par
  \vspace{0.45em}
}

\newcommand{\lignesreponse}[1]{%
  \begin{tcolorbox}[
    enhanced,
    colback=white,
    colframe=NeutreC!40,
    boxrule=0.5pt,
    arc=1.5mm,
    left=3mm,
    right=3mm,
    top=2mm,
    bottom=1mm,
    before skip=2pt,
    after skip=2pt
  ]
    \foreach \i in {1,...,#1}{\lignepointillee}
  \end{tcolorbox}
}

\newcommand{\imageactivite}[3]{%
  \IfFileExists{#3}{%
    \includegraphics[width=#1,height=#2,keepaspectratio]{#3}%
  }{%
    \fbox{%
      \parbox[c][#2][c]{#1}{%
        \centering\sffamily\small
        Image à intégrer\par
        \vspace{0.3em}
        \texttt{\detokenize{#3}}
      }%
    }%
  }%
}

\newcommand{\titreseq}[1]{%
\begin{center}
  {\sffamily\LARGE\bfseries\color{TitreC}
  Séquence 5 -- Statistiques et proportionnalité}\\[0.2em]
  {\sffamily\large\bfseries\color{SavoirC}#1}
\end{center}
}

\begin{document}

\titreseq{Cours -- Tableaux et diagrammes circulaires}

\section{Reconnaître une situation de proportionnalité}

\definition{Une \textbf{grandeur} est une quantité que l'on peut mesurer,
compter ou calculer. Une grandeur possède souvent une unité.}

\exemple{Dans une situation d'achat de pommes au poids, les grandeurs sont :

\[
\text{masse en kilogrammes}
\qquad\text{et}\qquad
\text{prix en euros}.
\]}

\definition{Deux grandeurs sont \textbf{proportionnelles} lorsque les valeurs
de l'une s'obtiennent en multipliant les valeurs de l'autre par un même nombre.}

\definition{Ce nombre est appelé \textbf{coefficient de proportionnalité}.}

\exemple{Si \(1\) kg de pommes coûte \(2{,}50\) €, alors :

\[
\text{prix}=\text{masse}\times2{,}50.
\]

Le coefficient permettant de passer de la masse au prix est \(2{,}50\).}

\erreur{Deux grandeurs qui augmentent ensemble ne sont pas nécessairement proportionnelles.

Par exemple, l'âge d'un enfant et l'âge de sa mère augmentent ensemble,
mais ne sont pas multipliés par un même nombre.}

\section{Tableaux de proportionnalité}

\definition{Un tableau est un \textbf{tableau de proportionnalité} lorsque
ses deux lignes contiennent les valeurs correspondantes de deux grandeurs proportionnelles.}

\exemple{
\[
\begin{array}{c|ccc}
\text{Masse en kg} & 1 & 3 & 5\\
\hline
\text{Prix en €} & 2{,}5 & 7{,}5 & 12{,}5
\end{array}
\]
}

\subsection{Méthodes pour compléter un tableau}

\methode{\textbf{Utiliser le coefficient de proportionnalité}

On multiplie toutes les valeurs d'une ligne par le même nombre pour obtenir l'autre ligne.}

\exemple{
\[
3\times2{,}5=7{,}5.
\]}

\methode{\textbf{Revenir à l'unité}

On cherche d'abord la valeur correspondant à \(1\), puis on calcule la valeur demandée.}

\exemple{Si \(4\) kg coûtent \(10\) € :

\[
10\div4=2{,}5.
\]

Un kilogramme coûte \(2{,}50\) €. Donc \(7\) kg coûtent :

\[
7\times2{,}5=17{,}5\text{ €}.
\]}

\methode{\textbf{Multiplier ou diviser une colonne}

Si une valeur d'une colonne est multipliée ou divisée par un nombre,
la valeur correspondante doit subir la même opération.}

\methode{\textbf{Utiliser l'additivité ou la soustraction}

On peut additionner ou soustraire deux colonnes complètes.}

\exemple{Si \(3\) objets coûtent \(6\) € et \(4\) objets coûtent \(8\) €,
alors \(7\) objets coûtent :

\[
6+8=14\text{ €}.
\]}

\remarque{Plusieurs méthodes peuvent être correctes.
On choisit celle qui est la plus simple avec les nombres du problème.}

\subsection{Calcul avec trois valeurs connues}

\propriete{Dans un tableau de proportionnalité :

\[
\begin{array}{c|cc}
 & a & b\\
\hline
 & c & x
\end{array}
\]

on a :

\[
a\times x=b\times c.
\]

Lorsque \(a\neq0\) :

\[
x=(b\times c)\div a.
\]}

\remarque{Cette propriété ne doit être utilisée que si la situation
est bien proportionnelle et si les grandeurs sont correctement placées dans le tableau.}

\section{Diagrammes circulaires}

\definition{Un \textbf{diagramme circulaire} représente une répartition par des secteurs.
L'angle de chaque secteur est proportionnel à l'effectif, à la fréquence
ou au pourcentage de la modalité.}

\propriete{Le cercle entier représente le total :

\[
360^\circ
\longleftrightarrow
\text{effectif total}
\longleftrightarrow
100\%.
\]}

\methode{Pour calculer l'angle d'un secteur à partir d'un effectif :

\[
\text{angle}
=
\text{effectif de la modalité}
\times
360
\div
\text{effectif total}.
\]}

\methode{Pour calculer l'angle à partir d'un pourcentage :

\[
\text{angle}
=
\text{pourcentage}
\times
360
\div
100.
\]}

\exemple{Pour une modalité représentant \(25\%\) du total :

\[
25\times360\div100=90.
\]

Le secteur mesure \(90^\circ\).}

\subsection{Construire un diagramme circulaire}

\methode{
\begin{enumerate}[label=\textbf{\arabic*.}]
  \item calculer l'angle de chaque secteur ;
  \item vérifier que la somme des angles vaut \(360^\circ\) ;
  \item tracer un cercle et un premier rayon ;
  \item reporter les angles successivement au rapporteur ;
  \item ajouter un titre, une légende, les unités, la source et l'effectif total.
\end{enumerate}}

\subsection{Lire un diagramme avec esprit critique}

\aretenir{Avant de faire confiance à un diagramme circulaire, on vérifie :

\begin{itemize}
  \item que les catégories forment bien une répartition complète ;
  \item que la somme des pourcentages vaut \(100\%\) ;
  \item que la somme des angles vaut \(360^\circ\) ;
  \item que la taille des secteurs correspond aux nombres annoncés ;
  \item que le titre, la source, la date et l'effectif étudié sont indiqués.
\end{itemize}}

\erreur{Un secteur plus grand doit représenter une proportion plus grande.
Une représentation en perspective ou en volume peut donner une impression trompeuse.}

\end{document}
